\documentclass{article}
\usepackage[left=2cm, right=2cm, top=2cm, bottom=2cm]{geometry}
\usepackage{graphicx}
\usepackage{amsmath}
\usepackage{amssymb}
\usepackage{amsthm}
\usepackage{fancyhdr}
\usepackage{verbatim}
\usepackage{listings}
\usepackage{xcolor}
\usepackage{pdfpages}
\usepackage{cancel}
\usepackage{physics}
\usepackage{esint}
\usepackage{braket}

\lstset{
    language=C++,
    basicstyle=\ttfamily\footnotesize,
    keywordstyle=\color{blue},
    commentstyle=\color{green},
    stringstyle=\color{red},
    numbers=left,
    numberstyle=\tiny\color{gray},
    frame=single,
    breaklines=true,
    captionpos=b,
}


\title{PHYS 427: Lecture \# 7}
\author{Cliff Sun}

\newtheorem{theorem}{Theorem}[section]
\newtheorem{lemma}[theorem]{Lemma}
\newtheorem{definition}[theorem]{Definition}
\newtheorem{conjecture}[theorem]{Conjecture}
\newtheorem{proposition}[theorem]{Proposition}
\newtheorem{corollary}[theorem]{Corollary}
\newtheorem{one minute paper}[theorem]{One Minute Paper}

\pagestyle{fancy}
\lhead{\textbf{\thepage}\ \ \nouppercase{\rightmark}}
\chead{PHYS 427: Lecture \# 7}
\rhead{Cliff Sun}

\begin{document}

\maketitle

\section*{Lecture Span}
\begin{enumerate}
    \item Helmholtz Free Energy
\end{enumerate}

\section*{Recap}

At high temperatures, each quadratic DOF gives $k_B T /2$ to the energy. 

\begin{theorem}
    Equipartition theorem. 
    \begin{equation}
        \epsilon = a_1 q_1^2 + a_2 q_2^2 + \dots + a_n q_n^2
    \end{equation}
    Here, $q_i$ are continuous variables varying from $-\infty$ to $\infty$. For example, 
    \begin{equation}
        \epsilon = \frac{p^2}{2m} + \frac{1}{2}m \omega^2 x^2
    \end{equation}
\end{theorem}

In the high temperature limit, you turn the partition function into integrals. Then 

\begin{equation}
    Z(T) \approx A\int dq_i \dots dq_n e^{-\epsilon_i(q_1,q_2,\dots,q_n)\beta}
\end{equation}

Then, 

\begin{align}
    p(\epsilon_i) &= e^{-\epsilon_i/k_B T}/Z(T) = A \frac{dq_1 \dots dq_N e^{-\beta \epsilon(q_i)}}{A \int [\dots]}
\end{align}

Then, 

\begin{equation}
    \langle \epsilon_{q_1} \rangle = \langle a_1 q_1^2 \rangle = \frac{\int dq_{i}\dots a_1q_1^2 e^{-\beta \epsilon(q_i)}}{Z(T)}
\end{equation}

The integrals upstairs and downstairs cancel. Then, this fraction is

\begin{equation}
    = \frac{\int_{-\infty}^{\infty} dq_1 a_1q_1^2 \exp(\dots a_1q_1)}{\int \exp(\dots a_1q_1)}
\end{equation}

\begin{align}
    &= -\frac{\partial}{\partial \beta}\ln \left[ \int dq_1 \exp(\dots a_1 q_1^2) \right] \\
    &= -\partial_{\beta}\ln \left( \beta^{-1/2} \right)\\
    &= \frac{1}{2\beta} = \frac{1}{2}k_B T
\end{align}

This is simply the equiparitition theorem (only valid in the high temperature limit!). 

\subsection*{Example: Maxwell velocity distribution}

Probability for an atom in a gas to have speed in the range $[v, v + dv]$. For this, use spherical coordinates. Here, 

\begin{equation}
    \langle v \rangle = \sqrt{\frac{8 k_B T}{\pi m}}
\end{equation}
\begin{equation}
    \langle v^2 \rangle = \frac{3k_B T}{m}
\end{equation}

\section*{Helmholtz Free Energy}

To recap, we have had the microcanonical ensemble ($U$ is fixed, and volume of the system is fixed), then the entropy of this system increases until maximum. 

The canonical ensemble is if the system is coupled to a reservoir $R$ that is much larger than the system. Both have a fixed temperature of $T = T_S = T_R$, fixed by the reservoir. Also, volume of the system is fixed. 

The total energy $U_R + U_S$ is fixed. But both can individually vary. Therefore, the entropy of the system is max because $R + S$ is closed.

Is there a quantity of the system alone that is at an extremum? This is the Helmholtz free energy! Here, the number of particles $N$ is fixed. 

Start with $S(U,V)$, then 

\begin{align}
    dS &= \partial_{U}S dU + \partial_V S dV \\
    &= \frac{1}{T}dU + \cancelto{0}{\frac{p}{T}dV}\\
    &= \frac{1}{T}dU
\end{align}

Assume that volume is fixed, or $dV = 0$. Here, $dS_{R + S} = 0 = dS_R + dS_S$. Then 

\begin{equation}
    = \frac{dU_R}{T} + dS_S
\end{equation}

Moreover, $dU_R = -dU_S$. Then, 
\begin{equation}
    dS = -\frac{dU_S}{T} + dS_S = -\frac{1}{T}d\left(U_S - TS_S\right)
\end{equation}
Thus, 
\begin{equation}
    -\frac{1}{T}dF = -\frac{1}{T}d\left( U_S - TS_S \right)
\end{equation}

Here, $F$ is the Helmholtz free energy. Therefore, if $R + S$ is at equilibrium, then $F$ is at a minimum. 

The thermodynamic identity is below:
\begin{equation}
    dS = \frac{1}{T}dU + \frac{p}{T}dV \iff dU = TdS - pdV
\end{equation}

This is also the first law of thermodynamics. $TdS$ is the heat and $pdV$ is the work done on the system. For Helmholtz free energy, 

\begin{align}
    dF &= dU - TdS - SdT \\
    &= TdS - pdV - TdS - SdT \\
    &= -pdV - SdT \\
    &= \frac{\partial F}{\partial V}dV + \frac{\partial F}{\partial T}dT
\end{align}

Therefore, $F = F(T,V)$. Canonical ensemble is fixed by the temperature and the volume. Therefore, 

\begin{equation}
    p = -\frac{\partial F}{\partial V}, \quad S = -\frac{\partial F}{\partial T}
\end{equation}

Is there a relationship between $F$ and $Z$? Consider a microcanonical ensemble, then 

\begin{equation}
    S_S = k_B \ln(\Omega)
\end{equation}

In the canonical ensemble, $Z$ is like $\Omega$ such that:

\begin{equation}
    F = -k_B T \ln(Z)
\end{equation}

A quick sketch of this formula:

\begin{align}
    F &= U - TS \\
    \iff S &= \frac{U - F}{T} \\
    \frac{\partial F}{\partial T} &= \frac{F - U}{T}
\end{align}

Here, equation (25) solves the above diff eq. Here, 

\begin{align}
    p &= k_B T \partial_{V}\ln (Z) \\
    S &= \partial_{T}(k_B T \ln(Z))
\end{align}

For an ideal gas, $Z = (n_Q V)^N/N!$. 

Quantum density:
\begin{equation}
    n_Q = \left( \frac{mk_B T}{2\pi \hbar^2} \right)^{3/2}
\end{equation}

Then we use the above equations to derive:

\begin{equation}
    p = \frac{ Nk_B T}{V} \iff pV = Nk_B T 
\end{equation}

Moreover, using Stirling's approximation:

\begin{align}
    S &= -F/T + U/T \\
    &= N k_B \ln(n_Q V) - Nk_B \ln N + N k_B + k_B T \partial_T \ln(T^{-3/2}) \\
    &= Nk_B \left[ \ln \left( \frac{n_Q V}{N} \right) + 5/2 \right]
\end{align}

This is called the Sackur-Tetrode equation. 

\subsection*{Example}

System with two energy levels. Calculate $S$ and $T \rightarrow \infty, T \rightarrow 0$. The partition function $e^0 + e^{-\beta \Delta} = 1 + e^{-\beta \Delta}$. Then, 

\begin{equation}
    S = \partial_T \left( k_B T \ln(Z) \right) = k_B T\ln\left( 1 + e^{-\Delta /k_B T} \right) + \frac{\Delta}{T}\frac{\exp(-\Delta/k_B T)}{1 + \exp(-\Delta/k_B T)}
\end{equation}

As $T \rightarrow 0$, $S \rightarrow 0$. As $T \rightarrow \infty$, then $S \rightarrow k_B \ln(2)$. 

\end{document}